\chapter{Implementation}
\label{ch:implementation}

This chapter bridges the conceptual approach presented in Chapter~\ref{ch:proposed_method} with the concrete implementation. Rather than detailing code line-by-line, we explain the design decisions made when translating concepts into working software: which technologies enable which capabilities, how we structured components for flexibility, which specific features and parameters we selected, and why.

\section{Technology Stack and Platform Choices}

\subsection{Python as Implementation Language}

We implement the entire system in Python 3.11, chosen specifically to support the reinforcement learning and evolutionary architecture search approach:

\begin{itemize}
    \item \textbf{Deep Learning Integration}: PyTorch, our chosen neural network framework, is primarily Python-based. Using Python throughout eliminates language barriers between geometric computation and neural network training.

    \item \textbf{Rapid Experimentation}: Reinforcement learning requires extensive experimentation---testing different reward functions, network architectures, and training strategies. Python's dynamic nature enables faster iteration than compiled languages.

    \item \textbf{Numerical Computing}: NumPy provides efficient array operations for EMS computation, heightmap updates, and collision detection without requiring low-level implementation.
\end{itemize}

\subsection{PyTorch for Neural Networks}

Among deep learning frameworks, we selected PyTorch 2.x for several implementation-critical reasons:

\begin{itemize}
    \item \textbf{Dynamic Computation Graphs}: Our action space varies from dozens to hundreds of candidates at each step. PyTorch's dynamic graphs handle variable-length inputs naturally, while static graph frameworks require complex padding strategies.

    \item \textbf{Debugging Transparency}: Reinforcement learning debugging requires inspecting Q-values, gradients, and attention weights during training. PyTorch's eager execution makes this straightforward.

    \item \textbf{Custom Architecture Support}: Set Transformers with induced attention blocks are not standard layers. PyTorch's modular design allows implementing custom attention mechanisms cleanly.
\end{itemize}

\section{Designing for Architectural Flexibility}

A core implementation principle is building flexibility into the system architecture. Since Chapter~\ref{ch:proposed_method} proposes using neurogenesis to discover optimal architectures, the implementation must support varying these architectural choices easily.

\subsection{Genome-Driven Configuration}

We design the neural network components to be fully configurable through a genome representation. While the Proposed Method conceptually describes using Set Transformers for attention, the implementation allows three attention options:

\begin{itemize}
    \item \textbf{Set Transformer}: With induced attention blocks using learnable inducing points
    \item \textbf{Standard Transformer}: Traditional self-attention without inducing points
    \item \textbf{No Attention}: Direct processing without cross-action reasoning
\end{itemize}

This design choice lets evolution decide whether attention is beneficial. The genome specifies \texttt{attention\_type} as a categorical gene that can mutate between these options. If attention hurts performance for our problem, evolution will discover architectures without it.

Similarly, we make evolvable:
\begin{itemize}
    \item Network depths (number of encoder layers)
    \item Layer widths (hidden dimensions)
    \item CNN configurations (number of channels, kernel sizes)
    \item Activation functions (ReLU, GELU, SiLU)
    \item Attention heads and inducing points
\end{itemize}

\subsection{Configuration Management}

All hyperparameters---both fixed training parameters and evolvable architectural choices---are managed through Python dataclasses. This provides:

\begin{itemize}
    \item \textbf{Type Safety}: Each parameter has an explicit type, catching errors at development time
    \item \textbf{Documentation}: Default values document baseline configurations
    \item \textbf{Serialization}: Configurations can be saved with experiment results for reproducibility
    \item \textbf{Genome Conversion}: The genome class has a \texttt{to\_config()} method that generates network configurations from genetic encodings
\end{itemize}

This separation between genome (evolvable choices) and configuration (complete system setup) allows evolution to explore architecture space while keeping training hyperparameters fixed.

\section{Feature Engineering: State and Action Representation}

The Proposed Method establishes that we need state features (global packing progress) and action features (candidate placement quality). Here we detail which specific features we selected and why.

\subsection{State Features: The 8-Dimensional Global Vector}

We represent the global state as 8 scalar features, chosen to capture complementary aspects of packing progress:

\begin{enumerate}
    \item \textbf{Packing Ratio} ($\in [0,1]$): Total volume placed divided by total bin volume. Measures overall efficiency.

    \item \textbf{Items Remaining Ratio} ($\in [0,1]$): Fraction of items not yet placed. Indicates progress through the sequence.

    \item \textbf{Largest EMS Dimensions} (3 features, $\in [0,1]$): Width, depth, and height of the largest remaining empty space, normalized by bin dimensions. Large well-shaped spaces indicate packing potential; small fragmented spaces suggest trouble ahead.

    \item \textbf{Average Heightmap} ($\in [0,1]$): Mean height across the bin's heightmap grid. Low values indicate flat packing; high values suggest vertical builds.

    \item \textbf{Weight Utilization} ($\in [0,1]$): Current weight divided by maximum capacity. Critical when weight constraints are active.

    \item \textbf{Bins Used Ratio} ($\in [0,1]$): Current bins used divided by maximum allowed. Penalizes using many bins.
\end{enumerate}

These 8 features balance conciseness (enabling fast neural network processing) with expressiveness (capturing key state properties). We intentionally avoid high-dimensional representations that would complicate learning.

\subsection{Action Features: The 25-Dimensional Placement Vector}

Each candidate action (bin, item, EMS, rotation) is encoded as 25 features. We organize these into conceptual groups:

\paragraph{Geometric Features (12 dimensions):}
\begin{itemize}
    \item Original item dimensions (3): The item's unrotated size
    \item Rotated dimensions (3): The item size after applying this action's rotation
    \item EMS corner position (3): Where the placement occurs
    \item EMS dimensions (3): The empty space size
\end{itemize}

These features let the network reason about spatial fit and positioning.

\paragraph{Fit Quality Features (4 dimensions):}
\begin{itemize}
    \item Slack space (3): Difference between EMS size and item size in each dimension. Small slack indicates tight fit.
    \item Tightness score (1): Count of dimensions where item perfectly fills the EMS. Higher values indicate less wasted space.
\end{itemize}

These quantify how well the item utilizes the chosen space.

\paragraph{Placement Context (9 dimensions):}
\begin{itemize}
    \item Volume utilization (1): Item volume divided by bin volume
    \item Container aspect ratios (2): Bin proportions affecting packing strategy
    \item Placement height (1): Vertical position in bin, affecting stability
    \item Local average height (1): Mean heightmap value in neighborhood, indicating local geometry
    \item Item weight ratio (1): Item weight relative to heaviest item
    \item Remaining weight capacity (1): Available weight budget
    \item Bin index (1): Which bin this placement targets
    \item Current bin utilization (1): How full the target bin is
\end{itemize}

These contextual features help the network evaluate placement consequences.

All features are normalized to $[0, 1]$ or $[-1, 1]$ ranges to facilitate neural network training.

\subsection{Heightmap Design Choices}

The Proposed Method introduces heightmaps for spatial reasoning. Here we detail how we configured this representation.

\paragraph{Grid Resolution:} We use adaptive resolution based on bin dimensions:
\begin{equation}
\text{resolution} = \max\left(1, \frac{\min(W, D)}{20}\right)
\end{equation}

This creates approximately 20 cells along the smaller bin dimension. For a $(100, 80, 60)$ bin, the resolution is 4 units, yielding a $(25 \times 20)$ grid. This balances spatial detail against computational cost---finer grids provide more precise geometry but increase memory and CNN computation.

\paragraph{Patch Size:} We extract $7 \times 7$ patches centered on each action's placement position. This size provides:
\begin{itemize}
    \item Sufficient context to capture local geometry (corners, valleys, walls)
    \item Small enough for efficient CNN processing across hundreds of actions
    \item Compatibility with standard CNN architectures (odd dimensions enable clear center definition)
\end{itemize}

Smaller patches (e.g., $3 \times 3$) lose context; larger patches (e.g., $15 \times 15$) dramatically increase computation with diminishing returns.

\paragraph{Normalization:} Heights are normalized by dividing by bin height $H$, mapping values to $[0, 1]$. This makes the representation scale-invariant across different bin sizes and improves neural network training by keeping inputs in a bounded range.

\paragraph{Boundary Handling:} When extracting patches near bin edges, out-of-bounds cells are set to zero (ground level). This naturally represents the bin boundary as open space.

\section{EMS-Based Geometry and Constraint Handling}

The Proposed Method uses Empty Maximal Spaces for collision-free action generation. The implementation follows the standard EMS approach~\cite{crainic2012packing} with adaptations for our constraints.

\subsection{EMS Management}

We implement the 6-way difference method: when an item is placed, each intersected EMS is split into up to 6 sub-spaces (left, right, front, back, bottom, top). We then prune dominated spaces (smaller EMS completely contained in larger ones) and limit the set to the 50 largest spaces by volume.

This limit prevents EMS proliferation while maintaining high-quality candidates. Since we sort by volume, we keep the most promising placement locations. Empirically, 50 spaces suffice even for complex instances---rarely do we need smaller EMS spaces after the top 50.

\subsection{Gravity Simulation}

EMS candidates may position items floating in mid-air. We apply gravity by finding the highest supporting surface: for each placed item with overlapping XY footprint, we track its top surface. The item drops to rest on the maximum such surface (or the ground if no overlap exists).

This ensures physical plausibility without explicit physics simulation. The $\mathcal{O}(n)$ complexity (checking against $n$ placed items) is acceptable for typical problem sizes (20-50 items).

\subsection{Constraint Validation}

The implementation validates four constraint types during action generation:

\begin{itemize}
    \item \textbf{Weight Capacity}: Track cumulative weight per bin, reject actions exceeding capacity
    \item \textbf{Incompatibility}: Maintain forbidden item pairs as a set, check before allowing placement
    \item \textbf{Affinity}: Track which items must stay together, reject placements that split affinity groups across bins
    \item \textbf{Relative Positioning}: After gravity application, verify that heavier items do not rest on lighter items by checking XY footprint overlaps
\end{itemize}

These checks filter the action space during enumeration. Only actions satisfying all constraints appear as valid options for the neural network to evaluate.

\section{Reward Function Design}

The Proposed Method introduces potential-based reward shaping. Here we detail how we operationalized this concept.

\subsection{Potential Function Components}

We design the potential function $\Phi(s)$ to balance two objectives:

\begin{equation}
\Phi(s) = \text{bin\_utilization} + 0.3 \cdot \text{ems\_quality}
\end{equation}

\paragraph{Bin Utilization:} We compute volume used divided by total bin volume, directly measuring packing efficiency.

\paragraph{EMS Quality:} Rather than using all EMS spaces, we focus on the top 3 largest spaces. We average their volumes and apply a cube root transformation:

\begin{equation}
\text{ems\_quality} = \frac{(\text{avg\_top\_3\_volume})^{1/3}}{(\text{bin\_volume})^{1/3}}
\end{equation}

The cube root provides diminishing returns---doubling space size doesn't double the benefit. This reflects that packing difficulty scales sub-linearly with space size.

The 0.3 weighting balances immediate efficiency (utilization) against future potential (space quality). We tuned this empirically: higher weights overemphasize space preservation at the expense of actually placing items; lower weights ignore fragmentation until it's too late.

\subsection{Shaped Reward Computation}

At each step, we compute:
\begin{equation}
r_{\text{shaped}} = \gamma \cdot \Phi(s') - \Phi(s)
\end{equation}

This provides continuous feedback about packing quality changes. Placing an item that improves utilization while maintaining good spaces yields positive reward. Fragmentation that reduces EMS quality yields negative reward.

\subsection{Terminal Rewards}

At episode end, we add bonuses and penalties:

\begin{itemize}
    \item \textbf{Completion Bonus} (+0.3): Reward successfully placing all items
    \item \textbf{Bin Efficiency Bonus} (+0.2 scaled): Reward using fewer bins, computed as $0.2 \cdot (1 - \text{bins\_used} / \text{max\_bins})$
    \item \textbf{Unplaced Item Penalty} (-0.2 scaled): Penalize failed placements, scaled by fraction of unplaced items
\end{itemize}

These terminal signals align the shaped rewards with the true objective (minimizing bins).

\section{Neural Network Architecture Design}

The Proposed Method establishes a multi-stream architecture combining state encoding, action encoding, heightmap processing, and attention. Here we explain how we structured these components.

\subsection{Component Architecture Choices}

\paragraph{State Encoder:} We use a simple MLP (Multi-Layer Perceptron) to process the 8-dimensional state vector. The baseline configuration uses 2 hidden layers of 256 units each with ReLU activations. This is deliberately lightweight---the state vector is low-dimensional and doesn't require deep processing.

\paragraph{Heightmap CNN:} We design a two-stage CNN:
\begin{itemize}
    \item \textbf{Convolutional Layers}: Two convolutional layers with channels $[16, 32]$ and $3 \times 3$ kernels with padding to preserve spatial dimensions. This extracts local spatial patterns.
    \item \textbf{Global Pooling}: Adaptive average pooling reduces each feature map to a single value, aggregating spatial information into a fixed-size vector regardless of input patch size.
    \item \textbf{Projection}: A final fully-connected layer projects to a 64-dimensional embedding.
\end{itemize}

We chose shallow CNNs (only 2 conv layers) because $7 \times 7$ patches are small---deeper networks would over-parameterize. The genome makes channel counts evolvable, allowing evolution to discover if deeper/wider CNNs improve performance.

\paragraph{Action Encoder:} The action encoder receives concatenated action features (25-dim) and heightmap embeddings (64-dim), totaling 89 dimensions. We process this through an MLP matching the state encoder architecture. This fusion combines geometric features with spatial context before attention processing.

\paragraph{Attention Mechanism:} As discussed in Section~\ref{sec:genome_flexibility}, we support three attention configurations:

\begin{itemize}
    \item \textbf{Set Transformer}: Uses Induced Set Attention Blocks (ISAB) with 32 learnable inducing points. This reduces complexity from $\mathcal{O}(N^2)$ to $\mathcal{O}(NM)$ where $N$ is actions and $M=32$ is inducing points.
    \item \textbf{Standard Transformer}: Uses PyTorch's built-in transformer encoder with 2 layers and 4 attention heads.
    \item \textbf{None}: Skips attention entirely, processing actions independently.
\end{itemize}

The genome's \texttt{attention\_type} gene selects among these. Evolution determines which provides the best performance-complexity tradeoff.

\paragraph{Q-Value Head:} The final component combines state and action embeddings. We broadcast the state embedding across all actions, concatenate with each action embedding, and pass through a 2-layer MLP that outputs a single Q-value per action. This design allows the network to modulate action values based on global state.

\subsection{Information Flow}

The complete forward pass follows this pipeline:

\begin{enumerate}
    \item State features $\rightarrow$ State Encoder $\rightarrow$ state embedding (256-dim)
    \item Heightmap patches $\rightarrow$ CNN $\rightarrow$ heightmap embeddings (64-dim per action)
    \item Action features + heightmap embeddings $\rightarrow$ Action Encoder $\rightarrow$ action embeddings (256-dim per action)
    \item Action embeddings $\rightarrow$ Attention (optional) $\rightarrow$ contextualized embeddings
    \item State embedding + action embeddings $\rightarrow$ Q-Head $\rightarrow$ Q-values
    \item Apply action mask to set invalid actions to $-\infty$
\end{enumerate}

This design separates concerns: spatial reasoning (CNN), global context (state encoder), geometric analysis (action encoder), action relationships (attention), and value estimation (Q-head).

\section{Training Configuration}

The Proposed Method uses Double DQN with experience replay. Here we detail the specific training choices.

\subsection{Experience Replay Design}

We use a circular buffer (implemented via Python's \texttt{collections.deque}) with capacity 200,000 transitions. Each stored transition includes:

\begin{itemize}
    \item State and next state vectors (8-dim each)
    \item Action index and reward
    \item Action features and heightmap patches for both current and next states
    \item Action validity masks indicating which actions are valid
    \item Episode termination flag
\end{itemize}

The large buffer capacity (200K) allows storing diverse experiences across many episodes. We set the warmup period to 1,000 steps, ensuring the buffer contains sufficient variety before training begins.

\subsection{Double DQN Implementation}

We implement Double DQN by maintaining two networks: an online network updated every training step, and a target network updated every 1,000 steps via hard copy. We chose hard updates over soft updates ($\tau$-weighted averaging) for stability---gradual target changes weren't necessary given our update frequency.

For target computation, the online network selects the best next action, while the target network evaluates that action's value. This decoupling reduces overestimation.

\subsection{Training Hyperparameters}

Key training decisions:

\begin{itemize}
    \item \textbf{Batch Size}: 128 samples per update, balancing gradient stability with computational efficiency
    \item \textbf{Learning Rate}: $1 \times 10^{-4}$ with Adam optimizer, chosen for stable convergence
    \item \textbf{Gradient Clipping}: Clip gradients to norm 1.0, preventing instability from large Q-value errors
    \item \textbf{Discount Factor}: $\gamma = 0.992$, emphasizing long-term planning over immediate rewards
    \item \textbf{N-Step Returns}: 15-step returns to accelerate credit assignment
    \item \textbf{Loss Function}: Smooth L1 (Huber) loss, more robust to outliers than MSE
\end{itemize}

\subsection{Exploration Strategy}

We use $\varepsilon$-greedy exploration with linear decay:
\begin{itemize}
    \item Start: $\varepsilon = 1.0$ (fully random)
    \item End: $\varepsilon = 0.15$ (15\% random actions)
    \item Decay: Linear over 20,000 steps
\end{itemize}

We maintain 15\% exploration even after decay (rather than decaying to 0.05 or 0.01) because the action space changes dramatically across problem states. Continued exploration helps discover better placements even late in training.

\subsection{Variable Action Space Handling}

A key implementation challenge is the variable number of valid actions (10-500+ depending on state). We handle this through:

\begin{itemize}
    \item \textbf{Dynamic Action Sets}: Generate valid actions on-the-fly each step
    \item \textbf{Padding and Masking}: Pad action feature matrices to a fixed size (512 actions), use boolean masks to indicate which are valid
    \item \textbf{Masked Q-Values}: Set invalid action Q-values to $-\infty$ before action selection
\end{itemize}

This allows batch processing on GPUs while respecting the varying action spaces across states.

\section{Evolutionary Architecture Search Implementation}
\label{sec:genome_flexibility}

The Proposed Method uses neurogenesis to discover optimal architectures. Here we detail how we encoded and evolved network configurations.

\subsection{Genome Design}

We represent each architecture as a genome with two types of genes:

\paragraph{Multiplicative Genes} (integer values that double or halve during mutation):
\begin{itemize}
    \item Hidden dimensions: 64, 128, 256, 512, 1024
    \item Encoder layers: 1, 2, 4, 8
    \item Head hidden size: 64, 128, 256, 512
    \item Attention heads: 1, 2, 4, 8
    \item Inducing points: 8, 16, 32, 64
    \item Patch size: 3, 5, 7, 11
\end{itemize}

This multiplicative mutation scheme enables large architectural changes while maintaining reasonable values. Multiplying or dividing by 2 creates discrete jumps in capacity.

\paragraph{Discrete Genes} (categorical choices that mutate randomly):
\begin{itemize}
    \item Attention type: \texttt{"standard"}, \texttt{"set\_transformer"}, \texttt{"none"}
    \item CNN channels: \texttt{[8, 16]}, \texttt{[16, 32]}, \texttt{[32, 64]}, \texttt{[16, 48]}
    \item Activation function: \texttt{"relu"}, \texttt{"gelu"}, \texttt{"silu"}
    \item Dropout rate: 0.0, 0.1, 0.2, 0.3
\end{itemize}

These categorical genes let evolution explore fundamentally different architectural patterns---e.g., whether attention helps at all.

\subsection{Fitness Function Design}

We evaluate each genome by:
\begin{enumerate}
    \item Building a DQN agent with the genome's configuration
    \item Training for 20 episodes on small training problems
    \item Evaluating on validation problems
    \item Computing a multi-objective fitness:
\end{enumerate}

\begin{equation}
\text{fitness} = 0.7 \cdot \text{utilization} + 0.2 \cdot (1 - \text{bins\_penalty}) + 0.1 \cdot (1 - \text{complexity\_penalty})
\end{equation}

The weighting emphasizes packing quality (70\%) while penalizing excessive bin usage (20\%) and network size (10\%). The complexity penalty prevents evolution from always choosing the largest networks.

\subsection{Evolutionary Process}

We implement a standard genetic algorithm:

\begin{itemize}
    \item \textbf{Population}: 20 genomes per generation
    \item \textbf{Elitism}: Top 2 genomes survive unchanged
    \item \textbf{Selection}: Tournament selection with size 3
    \item \textbf{Crossover}: Uniform crossover swapping genes between parents
    \item \textbf{Mutation}: Adaptive rate starting at 0.2, decaying to 0.05 over generations
    \item \textbf{Generations}: 10 generations (limited by computational budget)
\end{itemize}

The adaptive mutation rate allows broad exploration early (high mutation) while refining successful architectures later (low mutation).

\section{Summary}

This chapter detailed the implementation choices that translate the conceptual approach from Chapter~\ref{ch:proposed_method} into working software. The key design decisions include:

\paragraph{Technology Selection:} Python for rapid experimentation, PyTorch for dynamic neural networks, NumPy for geometric computation.

\paragraph{Feature Engineering:} 8 global state features capturing packing progress, 25 action features describing placement quality, and $7 \times 7$ heightmap patches for spatial context.

\paragraph{Architectural Flexibility:} Genome-driven configuration allowing evolution to choose between Set Transformer, standard attention, or no attention; variable network depths and widths; different activation functions.

\paragraph{Reward Shaping:} Potential function combining bin utilization with EMS quality (cube-root transformed, top-3 averaged, 0.3 weighted) plus terminal bonuses for completion and bin efficiency.

\paragraph{Training Setup:} 200K-transition replay buffer, Double DQN with hard target updates every 1,000 steps, 15-step returns, $\varepsilon$-greedy exploration (1.0 → 0.15 over 20K steps), batch size 128, learning rate $10^{-4}$.

\paragraph{Evolutionary Search:} Multiplicative genes (hidden dims, layers) and discrete genes (attention type, CNN channels), 20-genome population, tournament selection, adaptive mutation (0.2 → 0.05), multi-objective fitness (70\% utilization, 20\% bins, 10\% complexity).

These choices bridge the gap between what the system does conceptually (proposed method) and how it achieves it concretely (implementation). The next chapter evaluates this implementation experimentally, measuring packing performance and analyzing the impact of these design choices.
